\section{Conclusion}

Under the frozen public-data development benchmark, the tested QML candidates
did not outperform strong classical controls, and no tested candidate is
eligible for a Gate 6 mission experiment. The primary whole-state fidelity
kernel (Q01), projected-kernel follow-up (Q01b), feasibility-only kernel (FQK),
and all 16 bounded successors produced negative outcomes under their stated
criteria. The strongest successor appeared better than the C06 physics-residual
control, but its advantage over a like-for-like classical alternative was only
4.65\%, below the 5\% threshold fixed before execution.

The study's contribution is a reproducible negative result with an explicit
governance and human-centered claim structure. It reports what was tested, what
did not advance, what remains a future lesson, and what a human reviewer would
need to see before considering further authorization. Its practical societal
value is a more accountable way to evaluate AI evidence for high-consequence
decision contexts, not a deployed mission capability. It does not establish
quantum advantage, fuel savings, mission utility, hardware performance, flight
readiness, NASA approval, or an invented QML model. That restraint is necessary
for the result to be useful.
